\section{矩阵分析}
\begin{proposition}
%@see: https://www.cse.cuhk.edu.hk/~cslui/CSCI3320_LECTURES/matrix_differentiation.pdf
设\(\vb{f}\colon K^n \to K^m\)，
则\begin{equation*}
	\def\arraystretch{1.5}
	\pdv{\vb{f}}{\vb{x}}
	= \begin{bmatrix}
		\pdv{f_1}{x_1}
		& \pdv{f_1}{x_2}
		& \dots
		& \pdv{f_1}{x_n} \\
		\pdv{f_2}{x_1}
		& \pdv{f_2}{x_2}
		& \dots
		& \pdv{f_2}{x_n} \\
		\vdots & \vdots && \vdots \\
		\pdv{f_m}{x_1}
		& \pdv{f_m}{x_2}
		& \dots
		& \pdv{f_m}{x_n} \\
	\end{bmatrix}.
\end{equation*}
\end{proposition}

\begin{corollary}
%@see: https://www.cse.cuhk.edu.hk/~cslui/CSCI3320_LECTURES/matrix_differentiation.pdf
设\(\vb{f}\colon K \to K^m\)，
则\begin{equation*}
	\def\arraystretch{1.5}
	\pdv{\vb{f}}{x}
	= \begin{bmatrix}
		\pdv{f_1}{x} \\
		\pdv{f_2}{x} \\
		\vdots \\
		\pdv{f_m}{x} \\
	\end{bmatrix}.
\end{equation*}
\end{corollary}

\begin{corollary}
%@see: https://www.cse.cuhk.edu.hk/~cslui/CSCI3320_LECTURES/matrix_differentiation.pdf
设\(f\colon K^n \to K\)，
则\begin{equation*}
	\pdv{f}{\vb{x}}
	= \begin{bmatrix}
		\pdv{f}{x_1}
		& \pdv{f}{x_2}
		& \dots
		& \pdv{f}{x_n} \\
	\end{bmatrix}.
\end{equation*}
\end{corollary}

\begin{example}
计算\(\pdv{(\vb\alpha / \norm{\vb\alpha})}{\vb\alpha}\)，
其中\(\norm{\vb\alpha}\)是\(\vb\alpha\)的\(L_2\)范数.
\begin{solution}
设\(\vb\alpha = (a_1,\dotsc,a_n)^T\)，
则\begin{align*}
	\pdv{\norm{\vb\alpha}}{\vb\alpha}
	&= \pdv{\sqrt{a_1^2 + \dotsb + a_n^2}}{\vb\alpha} \\
	&= \begin{bmatrix}
		\pdv{\sqrt{a_1^2 + \dotsb + a_n^2}}{a_1}
		& \pdv{\sqrt{a_1^2 + \dotsb + a_n^2}}{a_2}
		& \dots
		& \pdv{\sqrt{a_1^2 + \dotsb + a_n^2}}{a_n}
	\end{bmatrix} \\
	&= \begin{bmatrix}
		\frac{a_1}{\sqrt{a_1^2 + \dotsb + a_n^2}}
		& \frac{a_2}{\sqrt{a_1^2 + \dotsb + a_n^2}}
		& \dots
		& \frac{a_n}{\sqrt{a_1^2 + \dotsb + a_n^2}}
	\end{bmatrix} \\
	&= \frac{1}{\sqrt{a_1^2 + \dotsb + a_n^2}}
		\begin{bmatrix}
			a_1 & a_2 & \dots & a_n
		\end{bmatrix}
	= \frac{\vb\alpha^T}{\norm{\vb\alpha}},
\end{align*}
因此\begin{align*}
	\pdv{(\vb\alpha / \norm{\vb\alpha})}{\vb\alpha}
	&= \frac{1}{\norm{\vb\alpha}^2}
		\left(
			\pdv{\vb\alpha}{\vb\alpha} \norm{\vb\alpha}
			- \vb\alpha \pdv{\norm{\vb\alpha}}{\vb\alpha}
		\right) \\
	&= \frac{1}{\norm{\vb\alpha}^2}
		\left(
			\vb{E}_3 \norm{\vb\alpha}
			- \vb\alpha \frac{\vb\alpha^T}{\norm{\vb\alpha}}
		\right) \\
	&= \frac{1}{\norm{\vb\alpha}}
		\left(
			\vb{E}_3
			- \frac{\vb\alpha \vb\alpha^T}{\norm{\vb\alpha}^2}
		\right).
\end{align*}
\end{solution}
\end{example}

\begin{proposition}
%@see: https://www.cse.cuhk.edu.hk/~cslui/CSCI3320_LECTURES/matrix_differentiation.pdf
设\(\vb{y} = \vb{A} \vb{x}\)，
其中\(\vb{A} \in K^{m \times n}, \vb{x} \in K^n, \vb{y} \in K^m\)，
且\(\vb{A}\)不依赖于\(\vb{x}\)，
则\begin{equation*}
	\pdv{\vb{y}}{\vb{x}}
	= \vb{A}.
\end{equation*}
\end{proposition}

\begin{proposition}
%@see: https://www.cse.cuhk.edu.hk/~cslui/CSCI3320_LECTURES/matrix_differentiation.pdf
设\(\vb{y} = \vb{A} \vb{f}(\vb{x})\)，
其中\(\vb{A} \in K^{m \times n}, \vb{x} \in K^s, \vb{y} \in K^m\)，
且\(\vb{A}\)不依赖于\(\vb{x}\)，
则\begin{equation*}
	\pdv{\vb{y}}{\vb{x}}
	= \vb{A} \pdv{\vb{f}}{\vb{x}}.
\end{equation*}
\end{proposition}

\begin{proposition}
%@see: https://www.cse.cuhk.edu.hk/~cslui/CSCI3320_LECTURES/matrix_differentiation.pdf
设\(z = \vb{y}^T \vb{A} \vb{x}\)，
其中\(\vb{A} \in K^{m \times n}, \vb{x} \in K^n, \vb{y} \in K^m, z \in K\)，
且\(\vb{A}\)不依赖于\(\vb{x}\)和\(\vb{y}\)，
则\begin{equation*}
	\pdv{z}{\vb{x}}
	= \vb{y}^T \vb{A},
	\qquad
	\pdv{z}{\vb{y}}
	= \vb{x}^T \vb{A}^T.
\end{equation*}
\end{proposition}

\begin{corollary}
%@see: https://www.cse.cuhk.edu.hk/~cslui/CSCI3320_LECTURES/matrix_differentiation.pdf
设\(z = \vb{x}^T \vb{A} \vb{x}\)，
其中\(\vb{A} \in K^{n \times n}, \vb{x} \in K^n, z \in K\)，
且\(\vb{A}\)不依赖于\(\vb{x}\)，
则\begin{equation*}
	\pdv{z}{\vb{x}}
	= \vb{x}^T (\vb{A} + \vb{A}^T).
\end{equation*}
\end{corollary}

\begin{corollary}
%@see: https://www.cse.cuhk.edu.hk/~cslui/CSCI3320_LECTURES/matrix_differentiation.pdf
设\(z = \vb{x}^T \vb{A} \vb{x}\)，
其中\(\vb{A} \in K^{n \times n}, \vb{x} \in K^n, z \in K\)，
\(\vb{A}\)是对称矩阵，
且\(\vb{A}\)不依赖于\(\vb{x}\)，
则\begin{equation*}
	\pdv{z}{\vb{x}}
	= 2 \vb{x}^T \vb{A}.
\end{equation*}
\end{corollary}

\begin{proposition}
%@see: https://www.cse.cuhk.edu.hk/~cslui/CSCI3320_LECTURES/matrix_differentiation.pdf
设\(w = \vb{y}^T \vb{x}\)，
其中\(w \in K\)，
而\(\vb{x},\vb{y} \in K^n\)都是\(\vb{z} \in K^m\)的函数，
则\begin{equation*}
	\pdv{w}{\vb{z}}
	= \vb{x}^T \pdv{\vb{y}}{\vb{z}}
	+ \vb{y}^T \pdv{\vb{x}}{\vb{z}}.
\end{equation*}
\end{proposition}

\begin{corollary}
%@see: https://www.cse.cuhk.edu.hk/~cslui/CSCI3320_LECTURES/matrix_differentiation.pdf
设\(w = \vb{x}^T \vb{x}\)，
其中\(w \in K\)，
而\(\vb{x} \in K^n\)是\(\vb{z} \in K^m\)的函数，
则\begin{equation*}
	\pdv{w}{\vb{z}}
	= 2 \vb{x}^T \pdv{\vb{x}}{\vb{z}}.
\end{equation*}
\end{corollary}

\begin{proposition}
%@see: https://www.cse.cuhk.edu.hk/~cslui/CSCI3320_LECTURES/matrix_differentiation.pdf
设\(w = \vb{y}^T \vb{A} \vb{x}\)，
其中\(\vb{A} \in K^{m \times n}, \vb{x} \in K^n, \vb{y} \in K^m, w \in K\)，
\(\vb{x},\vb{y}\)都是\(\vb{z} \in K^s\)的函数，
且\(\vb{A}\)不依赖于\(\vb{z}\)，
则\begin{equation*}
	\pdv{w}{\vb{z}}
	= \vb{x}^T \vb{A}^T \pdv{\vb{y}}{\vb{z}}
	+ \vb{y}^T \vb{A} \pdv{\vb{x}}{\vb{z}}.
\end{equation*}
\end{proposition}

\begin{corollary}
%@see: https://www.cse.cuhk.edu.hk/~cslui/CSCI3320_LECTURES/matrix_differentiation.pdf
设\(w = \vb{x}^T \vb{A} \vb{x}\)，
其中\(\vb{A} \in K^{m \times n}, \vb{x} \in K^n, w \in K\)，
\(\vb{x}\)是\(\vb{z} \in K^s\)的函数，
且\(\vb{A}\)不依赖于\(\vb{z}\)，
则\begin{equation*}
	\pdv{w}{\vb{z}}
	= \vb{x}^T (\vb{A}^T + \vb{A}) \pdv{\vb{x}}{\vb{z}}.
\end{equation*}
\end{corollary}

\begin{corollary}
%@see: https://www.cse.cuhk.edu.hk/~cslui/CSCI3320_LECTURES/matrix_differentiation.pdf
设\(w = \vb{x}^T \vb{A} \vb{x}\)，
其中\(\vb{A} \in K^{m \times n}, \vb{x} \in K^n, w \in K\)，
\(\vb{x}\)是\(\vb{z} \in K^s\)的函数，
\(\vb{A}\)是对称矩阵，
且\(\vb{A}\)不依赖于\(\vb{z}\)，
则\begin{equation*}
	\pdv{w}{\vb{z}}
	= 2 \vb{x}^T \vb{A} \pdv{\vb{x}}{\vb{z}}.
\end{equation*}
\end{corollary}

\begin{proposition}
%@see: https://www.cse.cuhk.edu.hk/~cslui/CSCI3320_LECTURES/matrix_differentiation.pdf
设\(\vb{A}(t) \in K^{m \times m}\)是关于\(t\)的非奇异的函数矩阵，
则\(\vb{A}\)的逆矩阵\(\vb{A}^{-1}\)对\(t\)的导数为\begin{equation*}
	\pdv{\vb{A}^{-1}}{t}
	= - \vb{A}^{-1} \pdv{\vb{A}}{t} \vb{A}^{-1}.
\end{equation*}
\end{proposition}

%TODO https://ccrma.stanford.edu/~dattorro/matrixcalc.pdf
